% Source: https://artofproblemsolving.com/community/c3h1469988_problem_from_aops_intro_to_cp
% Stable identity: aops-topic-1469988-q01
Let $O$ be the centre. The triangle fails to contain $O$ exactly when the three points lie on a common semicircle (then the triangle is on one side of the diameter); if they do not, every line through $O$ has vertices on both sides, and $O$ is inside. Cases with a point exactly opposite another have probability $0$.

For each point $P_i$, let $E_i$ be the event that the other two points lie on the half-circle that starts at $P_i$ and goes counterclockwise. Each of the other two points lands there with probability $\dfrac{1}{2}$, independently, so $P(E_i)=\dfrac{1}{4}$. If the three points lie on a semicircle, the first of them going counterclockwise has the other two within half a turn, so some $E_i$ happens; and two of these events cannot happen together (except with probability $0$). Hence
\[
P(\text{all three on a semicircle})=3\cdot\dfrac{1}{4}=\dfrac{3}{4},
\]
and the probability that the triangle contains the centre is $1-\dfrac{3}{4}=\boxed{\dfrac{1}{4}}$.

\begin{asy}
usepackage("amsmath");
size(7cm);
defaultpen(fontsize(10pt));
pair O=(0,0);
draw(circle(O,1));
pair A=dir(100),B=dir(220),C=dir(330);
draw(A--B--C--cycle,blue+linewidth(0.9));
dot(O,red); label("$O$",O,N,red);
dot(A);dot(B);dot(C);
pair O2=(2.8,0);
draw(circle(O2,1));
draw(arc(O2,1,-10,170),red+linewidth(1.6));
draw(O2+dir(-10)--O2+dir(170),gray+dashed);
pair A2=O2+dir(-10),B2=O2+dir(60),C2=O2+dir(125);
draw(A2--B2--C2--cycle,blue+linewidth(0.9));
dot(O2,red); label("$O$",O2+(0,-0.27),red);
dot(A2);dot(B2);dot(C2);
label("contains $O$",(0,-1.3)); label("all on a semicircle",(2.8,-1.3));
label("\mbox{}",(min(currentpicture,user=true)+max(currentpicture,user=true))/2,black); // keeps the text after the figure black when LaTeX typesets the labels inline
\end{asy}
